\section{Local convergence: polyhedral partitions and curved planar interfaces}\label{art:local}\label{art:positive}

Many composite objectives in nonsmooth optimization, including smooth losses
regularized by $\ell_1$, hinge, or group-norm penalties and maxima of smooth
convex functions, admit structured local representations only near their
minimizer $x_*$. Because an arbitrary initial matrix $H_0\succ0$ can send the
first BFGS step in almost any descent direction, a Euclidean ball around $x_*$
is not invariant under the iteration. Convexity and uniqueness of $x_*$ resolve
this immediately: sufficiently low sublevels $\{F\le F_*+\varepsilon\}$ lie
inside any prescribed neighborhood of $x_*$ and trap every Armijo run
independently of $H_0$.

We use this sublevel localization to establish two distinct mechanisms for local
convergence. First, in arbitrary dimension $\R^n$, finite polyhedral cell
partitions and local centered decompositions enjoy inward radial accuracy near
$x_*$. Second, in $\R^2$, we treat curved switching interfaces, along which inward
radial secant bounds fail, by constructing a moving contour-velocity comparison
field.

\subsection{Sublevel localization, polyhedral partitions, and composite penalties}

\begin{lemma}[Matrix-independent sublevel basin]\label{art:localization}
Let $F:\R^n\to\R$ be finite and convex with a unique minimizer $x_*$. Then $F$
is coercive, every sublevel of $F$ is compact, and for every neighborhood $U$
of $x_*$ there exists $\varepsilon>0$ such that
$\{x:F(x)\le F(x_*)+\varepsilon\}\subset U$. Consequently, every admissible run
starting in that sublevel remains in $U$ for all $k\ge0$, regardless of
$H_0\succ0$.
\end{lemma}
\begin{proof}
Choose $r>0$ with $\overline B_r(x_*)\subset U$. Continuity of $F$, compactness
of the sphere $\{\norm{u}=r\}$, and uniqueness of $x_*$ give
\[
 \delta=\min_{\norm{u}=r}\bigl(F(x_*+u)-F(x_*)\bigr)>0.
\]
For any $x$ with $\norm{x-x_*}\ge r$, evaluating convexity at the point at
distance $r$ along the segment from $x_*$ to $x$ yields
\[
 F(x)-F(x_*)\ge\frac{\norm{x-x_*}}r\delta.
\]
Thus $F$ is coercive with compact sublevels, $\varepsilon=\delta/2$ gives
$\{F\le F(x_*)+\varepsilon\}\subset B_r(x_*)\subset U$, and Armijo descent
preserves this sublevel.\qed
\end{proof}

Near $x_*$, a finite polyhedral partition simplifies dramatically: every closed
cell that does not contain $x_*$ stays bounded away from $x_*$, so within a
small ball around $x_*$ every active cell is a convex polyhedron containing
$x_*$ itself. Any inward ray $x-\tau(x-x_*)$ therefore remains inside a single
cell, giving quadratic inward accuracy even when gradients jump across cell
facets and without requiring BFGS to identify active manifolds.

\begin{theorem}[Local convergence on finite polyhedral partitions]\label{art:polyhedral}
Let $n\ge1$ and let $F:\R^n\to\R$ be finite and convex with a unique minimizer
$x_*$. Suppose an open neighborhood $U$ of $x_*$ is covered by finitely many
closed full-dimensional polyhedra $P_i$, and for each $i$ there exist an open
set $V_i\supset U\cap P_i$ and a function $f_i\in C^{1,1}_{\rm loc}(V_i)$
(not necessarily convex outside $U\cap P_i$) such that
\[
 F(x)=f_i(x)\qquad(x\in U\cap P_i).
\]
Then there exists $\varepsilon>0$, depending only on $F$, $U$, and the partition
data (and independent of $H_0$, $c_1$, and $c_2$), such that
every infinite admissible full-BFGS run with arbitrary $H_0\succ0$, fixed
$0<c_1<c_2<1$, and initial value $F(x_0)\le F(x_*)+\varepsilon$ satisfies
\[
 x_k\in U\quad(k\ge0),\qquad F(x_k)\longrightarrow F(x_*),\qquad x_k\longrightarrow x_*.
\]
\end{theorem}


\begin{proof}[Theorem~\ref{art:polyhedral}]
Each closed polyhedron $P_i$ not containing $x_*$ has an open complement
containing $x_*$; since there are finitely many cells, one ball around $x_*$
avoids all cells with $x_*\notin P_i$. Shrink to a ball
$B_\rho(x_*)\subset U\cap\bigcap_{i:x_*\in P_i}V_i$ on which the gradients
$\nabla f_i$ of all remaining extensions share a common Lipschitz constant $L$.

At any differentiability point $x\in B_\rho(x_*)$ of $F$, pick a cell $P_i$
containing $x$. Since $x_*\in P_i$ and $P_i$ is convex, the entire inward
segment $x-\tau(x-x_*)$ ($0\le\tau\le1$) lies in $U\cap P_i$, where $F=f_i$.
Differentiating along the ray at $\tau=0^+$ gives
\[
 \nabla F(x)^T(x-x_*)=\nabla f_i(x)^T(x-x_*),
\]
and the smooth Taylor bound for $f_i$ along the segment yields
\[
 0\le F(x-\tau r_x)-F(x)+\tau\nabla F(x)^Tr_x
       \le\tfrac12L\tau^2\norm{r_x}^2
       \le\tfrac12L\rho^2\tau^2,
\]
where the lower bound uses convexity of $F$. Lemma~\ref{art:localization}
supplies $\varepsilon>0$ with $K=\{F\le F(x_*)+\varepsilon\}\subset B_\rho(x_*)$,
so~\eqref{art:inward-condition} holds on $K$ with $M_\eta=L\rho^2$ and
$\theta_\eta=1/2$. Corollary~\ref{art:radial-convergence} and Armijo descent
complete the proof.\qed
\end{proof}

\begin{corollary}[Piecewise linear-quadratic and composite objectives]\label{art:plq}
Let $F:\R^n\to\R$ have a unique minimizer $x_*$, and suppose $F$ is either a
finite convex piecewise linear-quadratic function or the sum of a convex locally
$C^{1,1}$ loss and a finite convex piecewise linear-quadratic regularizer. Then
$F$ enjoys the local sublevel convergence guarantee of
Theorem~\ref{art:polyhedral}.
\end{corollary}
\begin{proof}
Here piecewise linear-quadratic means represented by quadratic polynomials on a
finite cover by closed full-dimensional polyhedra (possibly with distinct
Hessians on different cells). Each quadratic polynomial extends smoothly to
$\R^n$, and adding a locally $C^{1,1}$ loss preserves local $C^{1,1}$ cell
extensions, so Theorem~\ref{art:polyhedral} applies.\qed
\end{proof}

Localizing Theorem~\ref{art:centered} via Lemma~\ref{art:localization} yields a
parallel local guarantee that also encompasses \emph{nonpolyhedral} group-norm
regularizers whose kinks have curved cross-sections away from the origin.

\begin{corollary}[Local smooth-plus-homogeneous basin]\label{art:local-centered}
If the local representation $F(x)=\psi(x)+\gamma(x-x_*)$ of
Theorem~\ref{art:centered} holds on a neighborhood $U$ of a unique minimizer
$x_*$, then there exists $\varepsilon>0$, depending only on $F$ and $U$ (and
independent of $H_0$, $c_1$, and $c_2$), such that every admissible run with
$F(x_0)\le F(x_*)+\varepsilon$ stays in $U$ and satisfies $F(x_k)\to F(x_*)$
and $x_k\to x_*$.
\end{corollary}
\begin{proof}
Shrink $U$ to an open convex neighborhood $U_0\subseteq U$ of $x_*$ on which the
representation holds; Lemma~\ref{art:localization} supplies a compact sublevel
$\{F\le F(x_*)+\varepsilon\}\subset U_0$, and Theorem~\ref{art:centered}
applies.\qed
\end{proof}

To see both corollaries at work on composite regularization, consider first the
$\ell_1$-regularized least-squares (Lasso) objective
$F(x)=\tfrac12\norm{Ax-b}^2+\lambda\norm{x}_1$ ($A\in\R^{m\times n}$, $\lambda>0$)
with a unique minimizer $x_*$ (already covered by Corollary~\ref{art:plq}).
Setting \mbox{$I_0=\{i:x_{*,i}=0\}$} and $\sigma_i=\operatorname{sign}(x_{*,i})$
for $i\notin I_0$, on the open convex neighborhood
$\{x:\sigma_ix_i>0\text{ for all }i\notin I_0\}$ we have the exact centered
splitting
\begin{equation}\label{art:local-l1-decomposition}
 \begin{aligned}
 F(x)&=\psi(x)+\gamma(x-x_*),\\
 \psi(x)&=\tfrac12\norm{Ax-b}^2+\lambda\sum_{i\notin I_0}\sigma_ix_i,\qquad
 \gamma(r)=\lambda\sum_{i\in I_0}|r_i|,
 \end{aligned}
\end{equation}
placing it directly under Corollary~\ref{art:local-centered} as well. Now
replace $\lambda\norm{x}_1$ by a sum of Euclidean group-norm residuals
$\sum_j\norm{A_jx-b_j}_2$ added to any convex locally $C^{1,1}$ loss: near a
unique minimizer $x_*$, every group with $A_jx_*-b_j\ne0$ is smooth, while every
active group with $A_jx_*-b_j=0$ equals the centered homogeneous seminorm
$\norm{A_j(x-x_*)}_2$. Corollary~\ref{art:local-centered} therefore guarantees
local convergence even though $\norm{A_j(x-x_*)}_2$ is not polyhedral.

\subsection{Curved switching interfaces and contour-velocity fields in the plane}

What happens when nonsmooth kinks follow \emph{curved} switching manifolds
through $x_*$, as in a pointwise maximum $F=\max_{1\le i\le m}f_i$ of smooth
strongly convex functions whose active Hessians differ? In that setting, inward
rays toward $x_*$ cross the curved switching locus transversely, and the radial
secant condition~\eqref{art:radial-condition} can fail on \emph{every}
neighborhood of $x_*$. For instance, on the three-quadratic maximum of
Theorem~\ref{art:nonradial-example} below, the tie curve $f_0(u,v)=f_1(u,v)$
satisfies $v=u-\tfrac12u^2+O(u^3)$, along which the gradient jump pairs
nontrivially with the radial vector, $(\nabla f_0-\nabla f_1)^T(u,v)=\tfrac12u^2-v^2\ne0$,
driving $(y^Tz)^2/t\to+\infty$ for short steps straddling the curve.

To prove convergence across curved kinks, the comparison vector $w(x)$ must
rotate so that it pairs equally with both active gradients across each switching
curve while changing slowly enough along BFGS steps to satisfy the transport
bound of Lemma~\ref{art:tail-comparison}. In $\R^2$, we can build such a field globally
around any positive level curve away from the \emph{affine-critical set}
\[
 \mathcal C_{\rm aff}=\{p\in\R^2:0\in\aff\partial F(p)\}.
\]
Recall from Section~\ref{art:geometry} why $\mathcal C_{\rm aff}$ is the natural
critical set for equal-rate comparison directions: at any $p\notin\mathcal C_{\rm aff}$,
the minimum-norm vector $a\in\aff\partial F(p)$ is nonzero and satisfies
$g^Ta=\norm{a}^2>0$ for all $g\in\partial F(p)$, so $-a/\norm{a}^2$ decreases
every active supporting model at unit rate.

The geometric construction on a compact non-affine-critical subset of a level
curve $\{F=\ell\}$ ($\ell>F_*$) proceeds in three steps (proved in full in
Lemma~\ref{art:planar-field} of Appendix~\ref{art:planar-appendix}):
\begin{enumerate}
\item \emph{Corner velocities cancel jumps.} Where two convex $C^2$ branches
$f_1,f_2$ tie with linearly independent gradients, the implicit function theorem
parametrizes the corner curve $p_j(t)$ by the level $f_1(p_j(t))=f_2(p_j(t))=t$.
Differentiating in $t$ gives a velocity vector $v_j(t)=p_j'(t)$ satisfying
$\nabla f_1(p_j(t))^Tv_j(t)=\nabla f_2(p_j(t))^Tv_j(t)=1$: thus $v_j$ has
strictly positive pairing with the subdifferential while annihilating the
leading gradient jump $(\nabla f_2-\nabla f_1)^Tv_j=0$.
\item \emph{Normal-angle blending cancels tangent derivatives.} When positive
sublevels $\{F\le t\}$ are strictly convex (contain no boundary line segments),
distinct corners on $\{F=t\}$ have disjoint supporting-normal arcs. Writing
$n(\theta)=(\cos\theta,\sin\theta)^T$ and $T(\theta)=(-\sin\theta,\cos\theta)^T$,
we blend the positive corner pairings $n(\theta)^Tv_j(t)$ smoothly with $1$ into
a positive periodic height function $h(t,\theta)$ and set
$w(x)=h(F(x),\theta(x))n(\theta(x))+h_\theta(F(x),\theta(x))T(\theta(x))$ where
$n(\theta(x))=\nabla F(x)/\norm{\nabla F(x)}$. Near the $j$th switching curve,
$h=n^Tv_j$ and $h_\theta=T^Tv_j$, so $w(x)=v_j(F(x))$ extends Lipschitz-continuously
across the kink. Everywhere smooth, the classical support-parametrization
identity $\partial_\theta(hn+h_\theta T)=(h+h_{\theta\theta})T$~\cite{geom_Bogosel2023}
cancels the normal component of the angular derivative, so $Dw$ maps the
one-dimensional level-set tangent line into itself: $\det(s,Dw\,s)=0$ whenever
$s$ is tangent to $\{F=\text{const}\}$.
\item \emph{Charging field rotation to Armijo decrease.} Because weak Wolfe
keeps eligible steps uniformly transverse to $w$, $|\det(s,w)|\ge\kappa\norm{s}$.
Along an eligible step $s$, tangent preservation bounds the derivative of
$\log|\det(s,w)|$ by a multiple of the normal displacement, hence by the
variation of $F$ along the chord, which convexity bounds by $2d$. Expressing
$w(z)=\lambda w(x)+\beta s$ in the basis $(w(x),s)$ therefore yields
$\lambda=\det(s,w(z))/\det(s,w(x))$ with $|\log\lambda|\le C_1d$ and
$\lambda^2\le1+Ld$, charging all field rotation to the summable decreases $d_k$
without requiring $\sum_k\norm{s_k}<\infty$.
\end{enumerate}

\begin{theorem}[Affine-critical cluster points in the plane]\label{art:planar-cluster}
Let $F:\R^2\to\R$ be finite and convex, attain its minimum $F_*$, and have a
compact sublevel $K=\{x:F(x)\le b\}$ with $b>F_*$. Assume that:
\begin{enumerate}
\item for every $F_*<t\le b$, the boundary of $\{F\le t\}$ contains no
nontrivial line segment; and
\item at every $p\in K\setminus\mathcal C_{\rm aff}$, either $F$ is $C^2$ near
$p$, or $F=\max\{f_1,f_2\}$ near $p$ for convex $C^2$ functions $f_1,f_2$ both
active at $p$ with $\nabla f_1(p)$ and $\nabla f_2(p)$ linearly independent.
\end{enumerate}
Then every infinite admissible full-BFGS run starting in $K$ has at least one
cluster point in $\mathcal C_{\rm aff}$ (with no local structure assumed at
points of $\mathcal C_{\rm aff}$ itself).
\end{theorem}
\begin{proof}
Along the run, $x_k\in K$, $F(x_k)\downarrow\ell\ge F_*$, gradients are
uniformly bounded on $K$, and $\sum_kd_k<\infty$. Let $\Gamma\subset\{F=\ell\}$
be the nonempty compact cluster set of $\{x_k\}$. If $\ell=F_*$, then
$\Gamma\subseteq\argmin F\subseteq\mathcal C_{\rm aff}$. Suppose for
contradiction that $\Gamma\cap\mathcal C_{\rm aff}=\emptyset$; then $\ell>F_*$.

First, $\norm{s_k}\to0$. Otherwise compactness yields a subsequence with
$x_k\to x$, $x_{k+1}\to z\ne x$, and $d_k\to0$. For every $\theta\in[0,1]$,
\[
 F(x_k)-\theta d_k
 \le F((1-\theta)x_k+\theta x_{k+1})
 \le(1-\theta)F(x_k)+\theta F(x_{k+1}),
\]
so passing to the limit gives $F\equiv\ell$ on the segment $[x,z]$. Since
$\ell>F_*$, convexity forces every point of $\{F=\ell\}$ to lie on the boundary
of $\{F\le\ell\}$ (an interior point of $\{F\le\ell\}$ with value $\ell$ would
lie between a minimizer and a point of $\{F\le\ell\}$, forcing $F<\ell$),
contradicting hypothesis~1.

Applying Lemma~\ref{art:planar-field} (Appendix~\ref{art:planar-appendix}) to
the compact set $\Gamma\subset\{F=\ell\}\setminus\mathcal C_{\rm aff}$ supplies
an open neighborhood $U\supset\Gamma$, a vector field $w$, and constants
$\eta,\delta>0$, $C_0,L\ge0$. Since $\dist(x_k,\Gamma)\to0$, $\norm{s_k}\to0$,
and $d_k\to0$, there exists $J$ such that for all $k\ge J$, $x_k,x_{k+1}\in U$,
$\norm{s_k}<\delta$, and $d_k\le1$; Lemma~\ref{art:planar-field} then gives
$w_k=w(x_k)$ satisfying $(g_k^Tw_k)^2\ge\eta^2$, $(y_k^Tw_{k+1})^2/t_k\le C_0$,
and $w_{k+1}=\lambda_kw_k+\beta_ks_k$ with $\lambda_k^2\le1+Ld_k$, contradicting
Lemma~\ref{art:tail-comparison}.\qed
\end{proof}

\begin{corollary}[Planar convergence when affine criticality implies optimality]\label{art:planar-convergence}
Under the hypotheses of Theorem~\ref{art:planar-cluster}, if
$\mathcal C_{\rm aff}\cap K\subseteq\argmin F$, then every admissible run
starting in $K$ satisfies $F(x_k)\to F_*$ and $\dist(x_k,\argmin F)\to0$, with
$x_k\to x_*$ when the minimizer $x_*$ is unique.
\end{corollary}
\begin{proof}
Theorem~\ref{art:planar-cluster} provides a cluster point $p\in\mathcal C_{\rm aff}\cap K\subseteq\argmin F$.
Monotonicity of $F(x_k)$ and continuity of $F$ give $\lim_kF(x_k)=F(p)=F_*$, and
compactness of $K$ yields $\dist(x_k,\argmin F)\to0$ (hence $x_k\to x_*$ under
uniqueness).\qed
\end{proof}

\subsection{Nondegenerate planar smooth maxima and an explicit nonradial basin}

For a finite smooth maximum $F=\max_{1\le i\le m}f_i$ in $\R^2$, both the
isolation condition $\mathcal C_{\rm aff}=\{x_*\}$ near $x_*$ and two-branch
regularity away from $x_*$ follow automatically from standard multiplier
nondegeneracy at $x_*$, namely affine independence of the active gradients
together with strictly positive optimal convex weights.

\begin{corollary}[Nondegenerate planar smooth maxima]\label{art:planar-max}
Let $F:\R^2\to\R$ be finite and convex with a unique minimizer $x_*$, and
suppose $F=\max_{1\le i\le m}f_i$ near $x_*$ for convex $C^2$ branches $f_i$.
Assume that sufficiently low positive sublevels have no boundary line segments
(for instance, if the branches active at $x_*$ share a positive lower Hessian
bound near $x_*$), and that the active index set $I_*=\{i:f_i(x_*)=F(x_*)\}$ has
affinely independent gradients $\{\nabla f_i(x_*)\}_{i\in I_*}$ admitting
strictly positive weights
\[
 \lambda_i^*>0\quad(i\in I_*),\qquad \sum_{i\in I_*}\lambda_i^*=1,
 \qquad \sum_{i\in I_*}\lambda_i^*\nabla f_i(x_*)=0.
\]
Then there exists $\varepsilon>0$, independent of $H_0\succ0$ and of
$0<c_1<c_2<1$, such that every infinite admissible full-BFGS run with
$F(x_0)<F(x_*)+\varepsilon$ satisfies $F(x_k)\to F(x_*)$ and $x_k\to x_*$.
\end{corollary}
\begin{proof}
By continuity, restrict to a neighborhood of $x_*$ on which all $i\notin I_*$
are strictly inactive. Form the $3\times|I_*|$ matrix $A(x)$ with columns
$(\nabla f_i(x),1)^T$ ($i\in I_*$). Affine independence at $x_*$ means $A(x_*)$
has full column rank, which persists nearby, where any solution of
$A(p)\lambda=(0,0,1)^T$ is uniquely given by the continuous formula
$\lambda(p)=(A(p)^TA(p))^{-1}A(p)^T(0,0,1)^T$ with $\lambda(x_*)=\lambda^*>0$.
Shrink the neighborhood so that every coordinate of $\lambda(p)$ stays strictly
positive. Because $\partial F(p)=\conv\{\nabla f_i(p):i\text{ active at }p\}$
(each active gradient supports $F$ globally by convexity, while any vector
outside their convex hull is strictly separated from it, violating the
directional-derivative formula $F'(p;d)=\max_{i\text{ active}}\nabla f_i(p)^Td$),
any $p\in\mathcal C_{\rm aff}$ in this neighborhood has $0$ as an affine
combination of active gradients; extending by $0$ on inactive indices in $I_*$
forces that weight vector to equal $\lambda(p)>0$, so all $i\in I_*$ are active
and $0\in\conv\partial F(p)$, making $p=x_*$.

Moreover, at most three gradients in $\R^2$ can be affinely independent ($|I_*|\le3$).
At any $p\ne x_*$ in this neighborhood, $p\notin\mathcal C_{\rm aff}$, so three
branches cannot be simultaneously active (their affine hull would be all of
$\R^2$, containing $0$); one active branch makes $F$ locally $C^2$, and two
active branches must have linearly independent gradients (otherwise the line
through them would pass through $0$, putting $p\in\mathcal C_{\rm aff}$).
Applying Lemma~\ref{art:localization} and Corollary~\ref{art:planar-convergence}
gives $\varepsilon>0$. Finally, if the branches $i\in I_*$ satisfy $\nabla^2f_i\succeq\mu I\succ0$ on
a convex neighborhood of $x_*$, then $f_i(x)-\tfrac\mu2\norm{x}^2$ is convex for
each $i\in I_*$, so $F$ is $\mu$-strongly convex there and the boundaries of its
positive sublevels inside that neighborhood contain no nontrivial line
segments.\qed
\end{proof}

This multiplier nondegeneracy matches the active-set regularity conditions of
local bundle-Newton theory~\preprintcite[Propositions~5.5--5.6]{LW} and
Gebken~\cite{BG}, with a crucial difference: Corollary~\ref{art:planar-max}
proves convergence of unmodified BFGS from objective geometry alone, with no
assumed iterate or spectral asymptotics. We conclude the curved-interface theory
with an explicit three-quadratic maximum showing that moving comparison fields
strictly reach beyond radial analysis.

\begin{theorem}[Explicit convergence basin with curved, nonradial switching]\label{art:nonradial-example}
For the three-quadratic maximum on $\R^2$ given by
\[
 \begin{aligned}
 f_0(u,v)&=u+u^2+\tfrac12v^2,\qquad
 f_1(u,v)=v+\tfrac12u^2+\tfrac32v^2,\\
 f_2(u,v)&=-u-v+\tfrac32u^2+v^2,\qquad F(u,v)=\max_{0\le i\le2}f_i(u,v),
 \end{aligned}
\]
every infinite admissible full-BFGS run with $F(x_0)<1/100$, arbitrary
$H_0\succ0$, and any fixed $0<c_1<c_2<1$ satisfies $F(x_k)\to0$ and
$x_k\to(0,0)$. Nevertheless, the radial secant condition~\eqref{art:radial-condition}
fails on a positive sublevel band inside every neighborhood of $(0,0)$, even
when restricted to pairs eligible for $(c_1,c_2)=(1/4,1/2)$.
\end{theorem}
\begin{proof}
Each branch has Hessian between $I$ and $3I$, and $(f_0+f_1+f_2)/3=u^2+v^2$, so
$F(u,v)\ge u^2+v^2$, $(0,0)$ is the unique minimizer with $F_*=0$, and all
positive sublevels are strictly convex. At $(0,0)$, all three branches are
active with gradients $(1,0)^T$, $(0,1)^T$, $(-1,-1)^T$ and weights $1/3$.
To certify the explicit basin $\{F<1/100\}$, solve $A(u,v)\lambda=(0,0,1)^T$
explicitly: with $D=\det A(u,v)=3(1+v-uv)$,
\[
 D\lambda_0=1-4u+3v-7uv,\quad
 D\lambda_1=1+2u-3v-uv,\quad
 D\lambda_2=1+2u+3v+5uv.
\]
On the box $|u|,|v|\le1/10$, we have $D\ge267/100$ and the three numerators are
at least $23/100$, $49/100$, and $45/100$, so $\lambda_i(u,v)>0$ everywhere on
that box. Since $\{F\le b\}\subset\{u^2+v^2\le b\}\subset[-1/10,1/10]^2$ for any
$b\in(F(x_0),1/100)$, the argument of Corollary~\ref{art:planar-max} shows that
$\mathcal C_{\rm aff}\cap\{F\le b\}=\{(0,0)\}$ and two-branch regularity holds
throughout $\{F\le b\}\setminus\{(0,0)\}$; Corollary~\ref{art:planar-convergence}
gives $F(x_k)\to0$ and $x_k\to(0,0)$.

To exhibit the radial obstruction, fix small $u>0$ and let
\[
 v=v(u)=\frac{-1+\sqrt{1+4u+2u^2}}2=u-\tfrac12u^2+O(u^3)
\]
solve $f_0(u,v)=f_1(u,v)$ with $f_2(u,v)<f_0(u,v)$. Along this switching curve,
\[
 (\nabla f_0-\nabla f_1)^T(u,v)
 =(1+u)u-(1+2v)v=\tfrac12u^2-v^2=-\tfrac12u^2+O(u^3)\ne0
\]
for small $u>0$. Taking $x=(u,v+e)$ (uniquely on branch~$1$) and $z=(u,v-e)$
(uniquely on branch~$0$) as $e\downarrow0$, direct evaluation gives
\[
 d=2e(1+3v+3e),\qquad t=2e(1+2v+4e),\qquad \Delta=e(1+4v)+e^2.
\]
As $u,e\to0^+$, $\Delta/d\to1/2$ and $t/d\to1$, so for fixed small $u>0$ and all
sufficiently small $e>0$ the pair $(x,z)$ is eligible for $(c_1,c_2)=(1/4,1/2)$
and stays above half the positive tie value $f_0(u,v)>0$. Yet
$y^Tz\to\tfrac12u^2-v^2\ne0$ while $t\to0^+$, so $(y^Tz)^2/t\to+\infty$ on a
positive band inside any given neighborhood of $(0,0)$.\qed
\end{proof}

\subsection{Verifying the hypotheses and existence of admissible runs}

In applying Table~\ref{art:positive-outcomes} to a given convex objective $F$,
all algorithmic side conditions hold automatically along any admissible run on a
compact sublevel: weak Wolfe preserves $H_k\succ0$, Armijo gives
$\sum_kd_k<\infty$, and finite convexity bounds all subgradients on $K$.
Verification therefore reduces to a concrete geometric check on $F$:
\begin{itemize}
\item \emph{Smooth-plus-homogeneous route (Theorem~\ref{art:centered}):} exhibit
a decomposition $F(x)=\psi(x)+\gamma(x-x_*)$ around a minimizer $x_*$ on an open
convex domain $U$ containing $\{F\le F(x_0)\}$ (or locally around a unique
minimizer $x_*$ via Corollary~\ref{art:local-centered}).
\item \emph{Polyhedral route (Theorem~\ref{art:polyhedral}, Corollary~\ref{art:plq}):}
verify uniqueness of $x_*$ and locally $C^{1,1}$ cell formulas on a finite
polyhedral cover near $x_*$.
\item \emph{Planar smooth-maximum route (Corollary~\ref{art:planar-max}):} at
the unique minimizer $x_*$, check that the active gradients are affinely
independent with strictly positive convex multipliers $\lambda_i^*>0$, together
with a positive lower Hessian bound (or strict contours) on the active branches.
\end{itemize}

Appendix~\ref{art:hypothesis-examples} pairs these checks with explicit BFGS
trajectories and resolves when infinite admissible runs exist from arbitrary
initial data: Proposition~\ref{art:examples-mixed} writes out the band
certificate and a degenerating nonsmooth BFGS trajectory on $u^2+|v|$;
Proposition~\ref{art:examples-full-rank} gives a full-rank certificate in $\R^n$;
Lemma~\ref{art:examples-run-existence} proves that on any bounded-below finite
convex objective having no differentiable stationary point and at most countably
many nondifferentiable points along each line, infinite admissible runs exist
from \emph{every} differentiable $x_0$, $H_0\succ0$, and $0<c_1<c_2<1$,
establishing unrestricted run existence for a parameterized curved two-quadratic
family (Proposition~\ref{art:examples-two-quadratics}) and for the
three-quadratic objective of Theorem~\ref{art:nonradial-example}
(Corollary~\ref{art:examples-nonradial-runs}); finally,
Proposition~\ref{art:examples-no-step} exhibits a three-piece coercive cone and
initial ray on which no differentiable weak-Wolfe step exists, showing why
general nonsmooth convergence theorems quantify over admissible runs.
